\section{Introduction}
\label{sec:intro}

Asked what happens if you walk around a table, or whether a leaning stack of blocks will hold, people tend to picture the result rather than reason it out in words. Vision-language models (VLMs) usually do the opposite: they encode the image once and carry out every intermediate step in text. Several recent methods give models a visual intermediate state instead, by drawing sketches~\cite{hu2024sketchpad}, by producing latent visual tokens~\cite{yang2026mirage,li2025lvr}, or by supplying extra views~\cite{wang2026mindcube}. These methods change many things at once (training data, architecture, compute), so it is hard to tell whether the visual state itself is what helps.

We isolate that one factor. We build scenes in which the imagined state a question asks about can be \emph{rendered exactly} (\cref{fig:teaser}): the view from where the viewer would stand, the scene after two objects trade places, or the aftermath of releasing a stack of blocks. We then hold the model, the question and the prompt fixed and vary only the second image the model receives. This lets us ask three questions:
\begin{enumerate}[nosep,leftmargin=*]
  \item \textbf{How much would perfect imagination help?} An exact rendering of the imagined state (the oracle) is a reference point for methods that produce a visual intermediate state, under our rendering and prompt protocol.
  \item \textbf{Does the model actually use it?} Replacing the oracle with a misleading or an unrelated image tests whether the answer depends on what the visual state shows, rather than on having a second image at all.
  \item \textbf{How faithful must imagination be, and how close do models get on their own?} We degrade the oracle with controlled position noise, and we let the model build its own imagined state by writing down the scene and rendering it.
\end{enumerate}

Across four open VLMs, the answers to these questions differ by task, and in the 7--8B models the seen image competes with the imagined one (\cref{sec:results}). Code, data and the supplementary material (prompts, further controls, all significance tests) are at \url{https://github.com/Lkumar209/lvr-imagination}.
